% Fig. 1 -- Orbit-Evidence architecture. Included by paper/icc_main.tex via \input.
%
% ROUTING RULE, absolute: every connector is orthogonal. No to[out=..,in=..], no bend left/right,
% no splines, no line jumps. Curved violation arcs in the previous version dominated the figure and
% consumed the vertical whitespace the architecture needed; three straight rails carry the same
% information in a fraction of the height.
%
% BOX BORDERS. Every stage box is exactly 7mm tall (the two-line boxes overflow a 6.2mm minimum,
% the one-line box does not, so a single minimum makes them unequal). Half-height is 0.5185 units,
% so rails and observation connectors terminate at y = +-0.52: on the border, not 1mm inside it,
% where they rendered as dashed stubs floating in the box interior.
%
% GRID. Stage centres sit on a 3.8-unit pitch at x = 1.6, 5.4, 9.2, 13.0, 16.8, 20.6. Layer boxes
% sit at exactly the stage x they observe, so every observation connector is one vertical segment.
% Keep that invariant: if a layer moves, move it to another stage's x, never between two.
%
% RAIL SEPARATION. Three violation rails run right-to-left at distinct heights on distinct risers,
% so no rail crosses another:
%   L1.1  riser x=21.4  rail y=3.15  descends at x=1.6   (longest span, highest rail)
%   L1.4  riser x=19.8  rail y=2.25  descends at x=9.2
%   L3.1  riser x=16.8  rail y=1.35  descends at x=13.0  (shortest span, lowest rail)
% L1.3, the ordering-catchable contrast, runs below the pipeline at y=-0.95 on riser x=17.6, clear
% of the deployment observation connector at x=16.8.
%
% SIZING. No \resizebox: the unit vectors make the picture about \textwidth so declared font sizes
% reach the page unscaled. paper/scripts/check_glyphs.py fails the build below 6.4 pt for author-set
% text. If this figure needs room, move nodes -- do not scale it down.
\begin{figure*}[!t]
\centering
\begin{tikzpicture}[
  x=6.75mm, y=6.75mm,
  font=\small,
  >={Latex[length=1.8mm]},
  stage/.style={draw=black!65,line width=0.6pt,align=center,minimum height=7mm,
                text width=17.5mm,inner sep=1.5pt,fill=blue!4,font=\scriptsize},
  % Uniform 13.2mm: L1 wraps to four lines and the others to three, so without a minimum the
  % four frames had four different tops and the band did not read as a grid.
  layer/.style={draw=black!55,line width=0.6pt,align=left,text width=23mm,minimum height=13.2mm,
                inner sep=2.5pt,fill=black!3,font=\scriptsize},
  rel/.style={draw=black!60,line width=0.6pt,align=left,text width=36mm,
              inner sep=2.5pt,fill=blue!3,font=\scriptsize},
  dec/.style={draw=black!65,line width=0.6pt,align=center,text width=21mm,
              minimum height=4mm,inner sep=1pt,fill=white,font=\scriptsize},
  grp/.style={draw=black!30,line width=0.5pt,dash pattern=on 2pt off 1.5pt,inner sep=4pt},
  ok/.style={->,line width=0.7pt,black!70},
  obs/.style={->,line width=0.5pt,black!40},
  bus/.style={line width=0.5pt,black!40},
  no/.style={->,line width=0.7pt,red!60!black,dashed},
  noord/.style={->,line width=0.7pt,blue!50!black,densely dotted},
  tie/.style={<->,line width=0.6pt,black!55},
  lab/.style={font=\scriptsize,inner sep=1.2pt},
  % fill=white knocks the band border out from behind a legend: the contract and decision
  % headers sit on their dashed frames, which otherwise strike straight through the text.
  hdr/.style={font=\scriptsize\bfseries,text=black!75,anchor=west,fill=white,inner sep=1.5pt}]

% ---------------------------------------------------------------- the experiment, one baseline
% Band headers sit at x positions no connector occupies. At x=0 this one was struck through by the
% L1.1 riser descending at x=1.6, and the contract header by the L1 observation connector; both were
% invisible in the source and obvious in the render. If a connector moves, re-check these two.
\node[hdr] at (2.2,0.95) {Learning-assisted satellite experiment};
\node[stage] (pub) at ( 1.6,0) {published\\orbital state};
\node[stage] (vis) at ( 5.4,0) {predicted\\visibility};
\node[stage] (reg) at ( 9.2,0) {frozen\\registry};
\node[stage] (mod) at (13.0,0) {model, tracker,\\selection, gate};
\node[stage] (dep) at (16.8,0) {deployment};
\node[stage] (lab) at (20.6,0) {label\\closure};
\foreach \a/\b in {pub/vis, vis/reg, reg/mod, mod/dep, dep/lab} \draw[ok] (\a) -- (\b);

% decision and freeze markers, on the baseline and clear of every rail
\draw[black!55,line width=0.6pt] (7.3,-0.52) -- (7.3,0.52);
\node[font=\scriptsize,text=black!65,anchor=north] at (7.3,-0.55) {decision};
\draw[black!55,line width=0.6pt] (14.9,-0.52) -- (14.9,0.52);
\node[font=\scriptsize,text=black!65,anchor=north] at (14.9,-0.55) {freeze};

% ---------------------------------------------------------------- violation rails, orthogonal only
\draw[no] (21.4,0.52) -- (21.4,3.15) -- (1.6,3.15) -- (1.6,0.52);
\node[lab,text=red!60!black,anchor=south] at (11.5,3.20)
  {\rulename{L1.1} future publication $\rightarrow$ deployment feature};
\draw[no] (19.8,0.52) -- (19.8,2.25) -- (9.2,2.25) -- (9.2,0.52);
\node[lab,text=red!60!black,anchor=south] at (14.5,2.30)
  {\rulename{L1.4} later reference $\rightarrow$ frozen row membership};
\draw[no] (16.8,0.52) -- (16.8,1.35) -- (13.0,1.35) -- (13.0,0.52);
\node[lab,text=red!60!black,anchor=south] at (14.9,1.40)
  {\rulename{L3.1} post-freeze observation $\rightarrow$ model state};
\draw[noord] (20.6,-0.52) -- (20.6,-0.95) -- (17.6,-0.95) -- (17.6,-0.52);
\node[lab,text=blue!50!black,anchor=north] at (19.1,-1.00)
  {\rulename{L1.3} label used before closure};

% ---------------------------------------------------------------- the contract, gridded
\node[layer] (l1) at ( 1.6,-3.45) {\textbf{\rulename{L1}} availability and membership\\[0.5pt]
  publication clock, row membership, label closure};
\node[layer] (l2) at ( 5.4,-3.45) {\textbf{\rulename{L2}} physical and scheduling validity\\[0.5pt]
  predicted visibility, physical bounds};
\node[layer] (l3) at (13.0,-3.45) {\textbf{\rulename{L3}} model-state causality\\[0.5pt]
  state canaries, probe effectiveness, control};
\node[layer] (l4) at (16.8,-3.45) {\textbf{\rulename{L4}} independence and reproducibility\\[0.5pt]
  fold chronology, provenance, unit};
\node[grp,fit=(l1)(l2)(l3)(l4)] (contract) {};
% After the frame, so the white knockout wins. Drawn before it, the dashes struck through the text.
\node[hdr] at (6.4,-2.15) {Orbit-Evidence contract};

% One vertical observation connector per layer: each layer sits at the x of the stage it observes,
% so no connector bends and none crosses another. An earlier version routed a second L1 branch from
% the registry at x=9.2 leftwards along y=-1.65, which necessarily crossed the L2 connector at
% x=5.4 -- no orthogonal route between those two points avoids it. The registry stays visibly
% bound to L1 through the L1.4 rail, which descends at exactly x=9.2.
% Terminated on the node borders, not on a hand-computed y: box heights depend on how the text
% wraps, and a literal left the arrowheads floating short of some boxes and inside others.
\draw[obs] ( 1.6,-0.52) -- (l1.north);
\draw[obs] ( 5.4,-0.52) -- (l2.north);
\draw[obs] (13.0,-0.52) -- (l3.north);
\draw[obs] (16.8,-0.52) -- (l4.north);

% ---------------------------------------------------------------- relational validity, one baseline
\node[hdr] at (0,-5.45) {Relational validity};
\node[rel] (r1) at ( 3.6,-6.55) {\textbf{run A} $\leftrightarrow$ \textbf{run B}\\[0.5pt]
  provenance completeness (\rulename{L4.6}), registry invariance (\rulename{L1.4})};
\node[rel] (r2) at (11.6,-6.55) {\textbf{declared unit} $\leftrightarrow$ \textbf{coarser level}\\[0.5pt]
  $\mathrm{ICC}(1)$ against a permutation reference (\rulename{L4.7})};
\draw[tie] (r1.east) -- (r2.west);
% One vertical trunk from the contract band, then a horizontal bus into both relational boxes.
% The trunk alone previously ended between the two boxes, pointing at nothing.
\draw[bus] (7.6,-4.63) -- (7.6,-5.15) -- (3.6,-5.15);
\draw[bus] (7.6,-5.15) -- (11.6,-5.15);
\draw[obs] ( 3.6,-5.15) -- (r1.north);
\draw[obs] (11.6,-5.15) -- (r2.north);

% ---------------------------------------------------------------- decision, one horizontal arrow
\node[dec] (dp) at (19.9,-5.85) {PASS};
\node[dec] (dh) at (19.9,-6.55) {HALT};
\node[dec] (di) at (19.9,-7.25) {INDETERMINATE};
\node[grp,fit=(dp)(dh)(di)] (decbox) {};
\node[hdr,anchor=west] at (18.3,-5.28) {Decision};
\draw[ok] (r2.east) -- (decbox.west);
\end{tikzpicture}
\caption{Orbit-Evidence architecture. A learning-assisted satellite experiment acts from state
available at the decision instant $t_d$, freezes row and model state before deployment, and closes
retrospective labels later. Four executable validity layers observe it. Most rules inspect a single
execution; the relational rules compare executions or aggregation levels. Applicable checks resolve
to \textsc{pass} or \textsc{halt}, while statistically unsupported decisions return
\textsc{indeterminate} rather than \textsc{pass}. Red dashed rails are prohibited paths, each
\emph{chronologically consistent} so an ordering check cannot express it; the blue dotted rail is
one of the two classes ordering does catch.}
\label{fig:contract}
\end{figure*}
